\GalSourcePageBoundary{067}{L0066}{38}{38}

\begin{center}{\bfseries PROPOSITION I.}\end{center}

\noindent\textsc{Théorème.} --- {\itshape Soit une équation donnée, dont $a,b,c,\ldots$
sont les $m$ racines. Il y aura toujours un groupe de permutations
des lettres $a,b,c,\ldots$ qui jouira de la propriété suivante:\par

1\textsuperscript{o} Que toute fonction des racines, invariable\textsuperscript{(*)} par les substitutions
de ce groupe, soit rationnellement connue;\par

2\textsuperscript{o} Réciproquement, que toute fonction des racines, déterminable
rationnellement, soit invariable par les substitutions.\par}

(Dans le cas des équations algébriques, ce groupe n'est autre
chose que l'ensemble des $1.2.3\ldots m$ permutations possibles
sur les $m$ lettres, puisque, dans ce cas, les fonctions symétriques
sont seules déterminables rationnellement.)

(Dans le cas de l'équation $\dfrac{x^{n}-1}{x-1}=0$, si l'on suppose $a=r$,
$b=r^{g}$, $c=r^{g^{2}},\ldots$, $g$ étant une racine primitive, le groupe de
permutations sera simplement celui-ci:
\[
\begin{array}{c}
abcd\ldots k,\\
bcd\ldots ka,\\
cd\ldots kab,\\
\ldots\ldots\ldots\\
kabc\ldots i;
\end{array}
\]
dans ce cas particulier, le nombre des permutations est égal au
degré de l'équation, et la même chose aurait lieu dans les équations
dont toutes les racines seraient des fonctions rationnelles
les unes des autres.)

\emph{Démonstration.} --- Quelle que soit l'équation donnée, on pourra
trouver une fonction rationnelle $V$ des racines, telle que toutes

\GalHistoricalNote{\textsuperscript{(*)} Nous appelons ici invariable non seulement une fonction dont la forme est
invariable par les substitutions des racines entre elles, mais encore celle dont la
valeur numérique ne varierait pas par ces substitutions. Par exemple, si $Fx=0$
est une équation, $Fx$ est une fonction des racines qui ne varie par aucune permutation.

Quand nous disons qu'une fonction est rationnellement connue, nous voulons
dire que sa valeur numérique est exprimable en fonction rationnelle des coefficients
de l'équation et des quantités adjointes.}
